\chapter{警報信号}
% full version: chapters/14_pain.tex

\pencilfig{14}{\pencilart{14}}{痛みは火災報知器であって、温度計ではない。}

ほかの感覚はすべて、世界がどうなっているかを伝える。痛みは違う。痛みが伝えるのは、何かが体を脅かしているということであり、その音量の大部分は刺激ではなくマシン自身が決める。痛みは火災報知器であって、温度計ではない。

痛みのセンサは強いものにしか反応しない。40数度を超える熱、強い圧力、刺激性の物質である。触覚のセンサよりはるかに慣れにくい。脅威がある限り、報知器は鳴りやんではならないのだ。

\section*{2本の配線}

指をぶつけたときのことを思い出してほしい。まず鋭く正確な痛みが来る。どこで、いつ。一瞬遅れて、鈍く広がる長い痛みが来る。どれほどひどいか。これは2本の別々の配線である。一方は速く、その信号は数十ミリ秒で届く。もう一方は遅く、その信号は1秒ほどかかる。

\section*{ゲート}

ぶつけたところをなぜさするのか。脊髄の中で、痛みの信号は「ゲート」を通る。触覚は抑制性ニューロンを起こし、それがゲートを閉じかける。だから打ったところをさするのは本当に効く。ただし中程度の痛みに対してだけで、強い痛みはそれでも通ってしまう。強い痛みがそのとき自ら抑制性ニューロンをオフにするというのは、もともとのゲート理論の一部だが、直接の裏づけはまだない。

\pencilfig{126}{%
% gate: the pain wire drives the relay neuron; touch wakes an inhibitory neuron that closes the gate
\draw[pen=pcAmber] (0,1.35) -- (3.05,1.0); \draw[arr=pcAmber] (2.8,1.03) -- (3.08,1.0);
\node[lbl, text=pcAmber!70!black, anchor=east] at (-0.05,1.35) {痛み};
\draw[pen=pcBlue] (0,-0.15) -- (1.75,-0.15); \draw[arr=pcBlue] (1.5,-0.15) -- (1.78,-0.15);
\node[lbl, text=pcBlue, anchor=east] at (-0.05,-0.15) {触覚};
\draw[dotpen=pcRed, wash=pcRed] (2.05,-0.15) circle (0.26);
\draw[pcRed, line width=0.7pt, -{Bar[width=6pt]}] (2.25,0.05) -- (3.2,0.66);
\node[lbl, text=pcRed, anchor=north] at (2.05,-0.45) {ブレーキ};
\draw[dotpen=pcGray, wash=pcGray] (3.4,0.9) circle (0.34);
\node[lbl, anchor=south] at (3.4,1.3) {ゲート};
\draw[pen] (3.75,0.9) -- (5.6,0.9); \draw[arr] (5.35,0.9) -- (5.65,0.9);
\node[lbl, anchor=west] at (5.7,0.9) {脳へ};
}{脊髄のゲート。痛みの信号はゲート役のニューロンを通って脳へ向かい、触覚はゲートを閉じかける抑制性ニューロンを起こす。だから打ったところをさするのは効くが、中程度の痛みに対してだけである。}

\section*{上からの音量つまみ}

脳は自分で痛みを大きくも小さくもできる。上から、脳からゲートへ放送チャネルが下りてくる。\nt{SERO}、\nt{NORA}、そして脳が自ら分泌する、鎮痛薬に似た特別な物質である。危険な状況では、痛みはほとんど消えることがある。危険から逃げている者に足を引きずっている暇はない。楽になるだろうという期待も、一部は同じ物質を通して痛みを和らげる。

\section*{報知器が鳴りやまないとき}

痛みが繰り返されると感度が上がる。同じ打撃がもっと痛くなり始めるのだ。これも学習であり、傷が治るまで自分で自分を強めるループである。だが強いループには危険がある。傷がもう治っても「張りついて」鳴り続けることがあるのだ。第2章の、自らを維持し続けるニューロン群のように。

\section*{教師と割り込み}

痛みは教師でもある。ドーパミンの「予想よりよいか悪いか」にとって、痛みは負の報酬として働く。そして痛みは、ノルアドレナリンの連打のように、今の作業に割り込む。すべてを放り出して危険に対処せよ。最も速い反応、熱いものから手を引っこめる動作は、前の章のあの拮抗筋の組によって脊髄の中で直接閉じている。あなたが痛みを感じるより前に。

\fullversion{14}
